\documentclass[11pt]{article}

\usepackage[margin=1in]{geometry}
\usepackage{amsmath,amssymb,amsthm}
\usepackage{booktabs}
\usepackage{enumitem}
\usepackage{algorithm}
\usepackage{algpseudocode}
\usepackage{fancyhdr}
\usepackage{xcolor}
\usepackage[colorlinks=true,linkcolor=blue!50!black,urlcolor=blue!50!black]{hyperref}

\pagestyle{fancy}
\fancyhf{}
\lhead{MATH 4610 --- Numerical Analysis II}
\rhead{Homework 5}
\cfoot{\thepage}

\newcommand{\R}{\mathbb{R}}
\newcommand{\dd}{\,\mathrm{d}}
\newcommand{\Oh}{\mathcal{O}}

\theoremstyle{definition}
\newtheorem{problem}{Problem}
\newenvironment{solution}{\begin{proof}[Solution]}{\end{proof}}

\begin{document}

\begin{center}
  {\Large\bfseries Homework 5: Runge--Kutta Methods}\\[6pt]
  Daniel Okafor \quad\textbullet\quad MATH 4610, Fall 2026 \quad\textbullet\quad Due Wednesday, October 14
\end{center}

\noindent Throughout, we consider the initial value problem
\begin{equation}\label{eq:ivp}
  y'(t) = f\bigl(t, y(t)\bigr), \qquad y(t_0) = y_0,
\end{equation}
with $f$ smooth enough for every Taylor expansion we write down. An
$s$-stage explicit Runge--Kutta method with step size $h$ is
\begin{equation}\label{eq:rk}
  k_i = f\Bigl(t_n + c_i h,\; y_n + h\sum_{j=1}^{i-1} a_{ij} k_j\Bigr),
  \quad i = 1,\dots,s, \qquad
  y_{n+1} = y_n + h \sum_{i=1}^{s} b_i k_i,
\end{equation}
which we summarise by its Butcher tableau
\[
  \renewcommand{\arraystretch}{1.2}
  \begin{array}{c|c}
    \mathbf{c} & A \\ \hline
               & \mathbf{b}^{\top}
  \end{array}
\]
with $A = (a_{ij})$ strictly lower triangular. We always assume the row-sum
condition $c_i = \sum_j a_{ij}$.

%----------------------------------------------------------------------
\begin{problem}[Two-stage methods]\label{prob:two-stage}
Consider the general explicit two-stage method
\[
  k_1 = f(t_n, y_n), \quad
  k_2 = f(t_n + c_2 h,\; y_n + a_{21} h k_1), \quad
  y_{n+1} = y_n + h(b_1 k_1 + b_2 k_2).
\]
\begin{enumerate}[label=(\alph*)]
  \item Derive the conditions on $b_1, b_2, c_2, a_{21}$ for the method to have order two.
  \item Show that the midpoint method, Heun's method and Ralston's method all belong to this family, and give their tableaux.
  \item Show that no member of the family has order three.
\end{enumerate}
\end{problem}

\begin{solution}
(a) Write $f$, $f_t$, $f_y$ for $f$ and its partial derivatives evaluated at $(t_n, y_n)$, and suppose $y_n = y(t_n)$. Differentiating $y' = f(t,y)$ once gives $y'' = f_t + f_y f$, so the exact solution satisfies
\[
  y(t_n + h) = y_n + h f + \frac{h^2}{2}\bigl(f_t + f f_y\bigr) + \Oh(h^3).
\]
Expanding $k_2$ in two variables,
\[
  k_2 = f + c_2 h f_t + a_{21} h f f_y + \Oh(h^2),
\]
so
\[
  y_{n+1} = y_n + h (b_1 + b_2) f + h^2 b_2 \bigl(c_2 f_t + a_{21} f f_y\bigr) + \Oh(h^3).
\]
Matching coefficients of $h f$, $h^2 f_t$ and $h^2 f f_y$ gives the order-two conditions
\begin{equation}\label{eq:order2}
  b_1 + b_2 = 1, \qquad b_2 c_2 = \tfrac12, \qquad b_2 a_{21} = \tfrac12.
\end{equation}
The last two force $a_{21} = c_2$ (the row-sum condition), and then $b_2 = 1/(2c_2)$ and $b_1 = 1 - 1/(2c_2)$. So the order-two two-stage methods form a one-parameter family indexed by $c_2 \neq 0$.

(b) The three named methods correspond to $c_2 = \tfrac12$, $1$ and $\tfrac23$:
\[
  \renewcommand{\arraystretch}{1.25}
  \underbrace{\begin{array}{c|cc}
    0 & & \\
    \tfrac12 & \tfrac12 & \\ \hline
    & 0 & 1
  \end{array}}_{\text{midpoint}}
  \qquad\qquad
  \underbrace{\begin{array}{c|cc}
    0 & & \\
    1 & 1 & \\ \hline
    & \tfrac12 & \tfrac12
  \end{array}}_{\text{Heun}}
  \qquad\qquad
  \underbrace{\begin{array}{c|cc}
    0 & & \\
    \tfrac23 & \tfrac23 & \\ \hline
    & \tfrac14 & \tfrac34
  \end{array}}_{\text{Ralston}}
\]
Each satisfies~\eqref{eq:order2}; for example, for Ralston $b_2 c_2 = \tfrac34\cdot\tfrac23 = \tfrac12$.

(c) Apply any member of the family to the test equation $y' = y$, $y(0) = 1$. One step gives
\[
  y_1 = 1 + h(b_1 + b_2) + h^2 b_2 a_{21} = 1 + h + \frac{h^2}{2},
\]
whereas $y(h) = e^h = 1 + h + h^2/2 + h^3/6 + \Oh(h^4)$. The local error is $h^3/6 + \Oh(h^4)$ regardless of the free parameter, so the local error is never $\Oh(h^4)$ and the method is never of order three. (More generally, the elementary differential $f_y^2 f$ appears in $y'''$ with coefficient $1$, and a two-stage method has no way to produce it.)
\end{solution}

%----------------------------------------------------------------------
\begin{problem}[Classical RK4]\label{prob:rk4}
The classical fourth-order Runge--Kutta method has tableau
\[
  \renewcommand{\arraystretch}{1.25}
  \begin{array}{c|cccc}
    0        &          &          &   & \\
    \tfrac12 & \tfrac12 &          &   & \\
    \tfrac12 & 0        & \tfrac12 &   & \\
    1        & 0        & 0        & 1 & \\ \hline
             & \tfrac16 & \tfrac13 & \tfrac13 & \tfrac16
  \end{array}
\]
Verify that it satisfies all eight order conditions for order four.
\end{problem}

\begin{solution}
The eight conditions, indexed by the rooted trees with at most four vertices, are listed in Table~\ref{tab:order-conditions} together with their values for RK4. We used
\[
  \mathbf{c} = \bigl(0, \tfrac12, \tfrac12, 1\bigr)^{\!\top}, \qquad
  A\mathbf{c} = \bigl(0, 0, \tfrac14, \tfrac12\bigr)^{\!\top}, \qquad
  A\mathbf{c}^2 = \bigl(0, 0, \tfrac18, \tfrac14\bigr)^{\!\top}, \qquad
  A^2\mathbf{c} = \bigl(0, 0, 0, \tfrac14\bigr)^{\!\top},
\]
where $\mathbf{c}^2$ means the componentwise square.

\begin{table}[ht]
  \centering
  \caption{Order conditions up to order four, evaluated for classical RK4.}
  \label{tab:order-conditions}
  \begin{tabular}{@{}clcc@{}}
    \toprule
    Order & Condition & Required & RK4 value \\
    \midrule
    1 & $\sum_i b_i$                      & $1$            & $\tfrac16 + \tfrac13 + \tfrac13 + \tfrac16 = 1$ \\[2pt]
    2 & $\sum_i b_i c_i$                  & $\tfrac12$     & $\tfrac16 + \tfrac16 + \tfrac16 = \tfrac12$ \\[2pt]
    3 & $\sum_i b_i c_i^2$                & $\tfrac13$     & $\tfrac1{12} + \tfrac1{12} + \tfrac16 = \tfrac13$ \\
      & $\sum_{i,j} b_i a_{ij} c_j$       & $\tfrac16$     & $\tfrac13\cdot\tfrac14 + \tfrac16\cdot\tfrac12 = \tfrac16$ \\[2pt]
    4 & $\sum_i b_i c_i^3$                & $\tfrac14$     & $\tfrac1{24} + \tfrac1{24} + \tfrac16 = \tfrac14$ \\
      & $\sum_{i,j} b_i c_i a_{ij} c_j$   & $\tfrac18$     & $\tfrac13\cdot\tfrac12\cdot\tfrac14 + \tfrac16\cdot 1\cdot\tfrac12 = \tfrac18$ \\
      & $\sum_{i,j} b_i a_{ij} c_j^2$     & $\tfrac1{12}$  & $\tfrac13\cdot\tfrac18 + \tfrac16\cdot\tfrac14 = \tfrac1{12}$ \\
      & $\sum_{i,j,k} b_i a_{ij} a_{jk} c_k$ & $\tfrac1{24}$ & $\tfrac16\cdot\tfrac14 = \tfrac1{24}$ \\
    \bottomrule
  \end{tabular}
\end{table}

All eight hold, so RK4 has order (at least) four. It does not have order five: applied to $y' = \lambda y$ it reproduces the Taylor series of $e^{h\lambda}$ only through the $z^4/24$ term (see Problem~\ref{prob:stability}).
\end{solution}

%----------------------------------------------------------------------
\begin{problem}[Absolute stability]\label{prob:stability}
\begin{enumerate}[label=(\alph*)]
  \item Find the stability function $R(z)$ of RK4, where $y_{n+1} = R(h\lambda) y_n$ when the method is applied to $y' = \lambda y$.
  \item Determine the interval of absolute stability on the negative real axis.
  \item For $y' = -50y$, what is the largest step size for which RK4 is absolutely stable? Compare with forward Euler.
\end{enumerate}
\end{problem}

\begin{solution}
(a) With $f(t,y) = \lambda y$ and $z = h\lambda$, the stages are
\begin{align*}
  h k_1 &= z y_n, \\
  h k_2 &= z\bigl(1 + \tfrac{z}{2}\bigr) y_n, \\
  h k_3 &= z\bigl(1 + \tfrac{z}{2} + \tfrac{z^2}{4}\bigr) y_n, \\
  h k_4 &= z\bigl(1 + z + \tfrac{z^2}{2} + \tfrac{z^3}{4}\bigr) y_n,
\end{align*}
and combining them with weights $\tfrac16, \tfrac13, \tfrac13, \tfrac16$ gives
\[
  R(z) = 1 + z + \frac{z^2}{2} + \frac{z^3}{6} + \frac{z^4}{24}.
\]
This is the degree-four Taylor polynomial of $e^z$, as it must be for any four-stage method of order four.

(b) For real $z<0$ we need $|R(z)| \le 1$. Since $R(z) > 0$ for all real $z$ (its minimum on $\R$ is about $0.27$, attained near $z \approx -1.6$), the binding condition is $R(z) \le 1$, i.e.
\[
  z\Bigl(1 + \frac{z}{2} + \frac{z^2}{6} + \frac{z^3}{24}\Bigr) \le 0.
\]
For $z<0$ this requires the cubic in parentheses to be non-negative. Its only real root, found by Newton's method from $z_0 = -3$, is $z^\ast \approx -2.7853$. The interval of absolute stability on the real axis is therefore
\[
  -2.7853 \le h\lambda \le 0.
\]

(c) With $\lambda = -50$ we need $50h \le 2.7853$, i.e.\ $h \le 0.0557$. Forward Euler has $R(z) = 1 + z$ and is stable for $-2 \le z \le 0$, so it needs $h \le 0.04$. RK4 allows a step about 39\% larger, but each step costs four function evaluations instead of one, so for a problem this stiff neither method is attractive and an implicit method should be used.
\end{solution}

%----------------------------------------------------------------------
\begin{problem}[Implementation and convergence]\label{prob:impl}
Implement forward Euler, Heun's method and RK4 and apply them to
\begin{equation}\label{eq:test}
  y' = y - t^2 + 1, \qquad 0 \le t \le 2, \qquad y(0) = 0.5,
\end{equation}
whose exact solution is $y(t) = (t+1)^2 - \tfrac12 e^{t}$. Report the error at $t = 2$ for $h = 0.2, 0.1, 0.05, 0.025$ and estimate the observed order of convergence.
\end{problem}

\begin{solution}
Algorithm~\ref{alg:rk4} gives the RK4 integrator in the vector-valued form we implemented (the scalar case is $m = 1$). Euler and Heun were implemented the same way with the obvious changes to the stage computations.

\begin{algorithm}[ht]
  \caption{Classical RK4 with fixed step size}
  \label{alg:rk4}
  \begin{algorithmic}[1]
    \Require $f:\R\times\R^m\to\R^m$, initial time $t_0$, final time $T$, initial value $\mathbf{y}_0\in\R^m$, number of steps $N$
    \Ensure Approximations $\mathbf{w}_n \approx \mathbf{y}(t_n)$ for $n = 0,\dots,N$
    \State $h \gets (T - t_0)/N$
    \State $t \gets t_0$, \quad $\mathbf{w}_0 \gets \mathbf{y}_0$
    \For{$n = 0, 1, \dots, N-1$}
      \State $\mathbf{k}_1 \gets f(t,\; \mathbf{w}_n)$
      \State $\mathbf{k}_2 \gets f(t + h/2,\; \mathbf{w}_n + \tfrac{h}{2}\mathbf{k}_1)$
      \State $\mathbf{k}_3 \gets f(t + h/2,\; \mathbf{w}_n + \tfrac{h}{2}\mathbf{k}_2)$
      \State $\mathbf{k}_4 \gets f(t + h,\; \mathbf{w}_n + h\,\mathbf{k}_3)$
      \State $\mathbf{w}_{n+1} \gets \mathbf{w}_n + \tfrac{h}{6}\bigl(\mathbf{k}_1 + 2\mathbf{k}_2 + 2\mathbf{k}_3 + \mathbf{k}_4\bigr)$
      \State $t \gets t_0 + (n+1)h$ \Comment{avoid accumulating round-off in $t$}
    \EndFor
    \State \Return $(\mathbf{w}_0, \dots, \mathbf{w}_N)$
  \end{algorithmic}
\end{algorithm}

The exact value is $y(2) = 9 - \tfrac12 e^2 \approx 5.3054719505$. Table~\ref{tab:errors} lists the global errors $|w_N - y(2)|$. The observed order between consecutive rows is
\[
  p_{\text{obs}} = \log_2 \frac{E(h)}{E(h/2)}.
\]

\begin{table}[ht]
  \centering
  \caption{Global error at $t = 2$ for problem~\eqref{eq:test}, and observed order of convergence.}
  \label{tab:errors}
  \begin{tabular}{@{}r rr rr rr@{}}
    \toprule
    & \multicolumn{2}{c}{Euler} & \multicolumn{2}{c}{Heun} & \multicolumn{2}{c}{RK4} \\
    \cmidrule(lr){2-3}\cmidrule(lr){4-5}\cmidrule(l){6-7}
    $h$ & Error & $p_{\text{obs}}$ & Error & $p_{\text{obs}}$ & Error & $p_{\text{obs}}$ \\
    \midrule
    0.200 & $4.397\times10^{-1}$ & ---  & $7.242\times10^{-2}$ & ---  & $1.089\times10^{-4}$ & ---  \\
    0.100 & $2.420\times10^{-1}$ & 0.86 & $1.890\times10^{-2}$ & 1.94 & $6.990\times10^{-6}$ & 3.96 \\
    0.050 & $1.275\times10^{-1}$ & 0.92 & $4.820\times10^{-3}$ & 1.97 & $4.421\times10^{-7}$ & 3.98 \\
    0.025 & $6.550\times10^{-2}$ & 0.96 & $1.216\times10^{-3}$ & 1.99 & $2.779\times10^{-8}$ & 3.99 \\
    \bottomrule
  \end{tabular}
\end{table}

The observed orders approach $1$, $2$ and $4$ respectively, as theory predicts. The efficiency gap is large: RK4 with $h = 0.2$ (40 function evaluations) is more accurate than Heun with $h = 0.025$ (160 evaluations) by an order of magnitude, and more accurate than Euler with $h = 0.025$ (80 evaluations) by nearly three orders.

For reference, Table~\ref{tab:traj} shows the RK4 trajectory with $h = 0.2$. The error grows steadily along the trajectory, which is expected because $\partial f/\partial y = 1 > 0$ amplifies earlier local errors by a factor of about $e^{h}$ per step.

\begin{table}[ht]
  \centering
  \caption{RK4 with $h = 0.2$ on problem~\eqref{eq:test}.}
  \label{tab:traj}
  \begin{tabular}{@{}cccc@{}}
    \toprule
    $t_n$ & $w_n$ & $y(t_n)$ & $|w_n - y(t_n)|$ \\
    \midrule
    0.0 & 0.5000000 & 0.5000000 & 0 \\
    0.4 & 1.2140762 & 1.2140877 & $1.14\times10^{-5}$ \\
    0.8 & 2.1272027 & 2.1272295 & $2.69\times10^{-5}$ \\
    1.2 & 3.1798942 & 3.1799415 & $4.74\times10^{-5}$ \\
    1.6 & 4.2834095 & 4.2834838 & $7.43\times10^{-5}$ \\
    2.0 & 5.3053630 & 5.3054720 & $1.09\times10^{-4}$ \\
    \bottomrule
  \end{tabular}
\end{table}
\end{solution}

%----------------------------------------------------------------------
\begin{problem}[Adaptive step size]\label{prob:adaptive}
Embedded pairs compute two solutions of different orders from the same stages and use their difference as an error estimate.
\begin{enumerate}[label=(\alph*)]
  \item Write pseudocode for an adaptive integrator based on the Runge--Kutta--Fehlberg 4(5) pair, using the step-size update $h_{\text{new}} = q\,h$ with $q = 0.84\,(\varepsilon/\hat{E})^{1/4}$.
  \item Explain where the exponent $1/4$ comes from.
  \item Run your code on~\eqref{eq:test} with tolerances $\varepsilon = 10^{-4}, 10^{-6}, 10^{-8}$ and report the number of accepted and rejected steps and the actual error at $t = 2$.
\end{enumerate}
\end{problem}

\begin{solution}
(a) Algorithm~\ref{alg:rkf45} is the controller we used. Here $\mathbf{w}^{(4)}$ and $\mathbf{w}^{(5)}$ are the fourth- and fifth-order results built from the six Fehlberg stages; we propagate the fourth-order value, as in Fehlberg's original method. The error estimate $\hat{E}$ is per unit step, so the tolerance controls the error committed per unit of $t$.

\begin{algorithm}[ht]
  \caption{Adaptive RKF45 integration}
  \label{alg:rkf45}
  \begin{algorithmic}[1]
    \Require $f$, $t_0$, $T$, $\mathbf{y}_0$, tolerance $\varepsilon$, initial step $h$, bounds $0 < q_{\min} < 1 < q_{\max}$
    \State $t \gets t_0$, \quad $\mathbf{w} \gets \mathbf{y}_0$
    \While{$t < T$}
      \State $h \gets \min(h,\; T - t)$ \Comment{do not step past $T$}
      \State compute stages $\mathbf{k}_1,\dots,\mathbf{k}_6$ from $(t, \mathbf{w}, h)$
      \State $\mathbf{w}^{(4)} \gets \mathbf{w} + \sum_{i} b_i \mathbf{k}_i$, \quad $\mathbf{w}^{(5)} \gets \mathbf{w} + \sum_i \hat{b}_i \mathbf{k}_i$
      \State $\hat{E} \gets \lVert \mathbf{w}^{(5)} - \mathbf{w}^{(4)} \rVert_\infty / h$
      \If{$\hat{E} \le \varepsilon$}
        \State $t \gets t + h$, \quad $\mathbf{w} \gets \mathbf{w}^{(4)}$ \Comment{accept}
        \State record $(t, \mathbf{w})$
      \EndIf
      \State $q \gets 0.84\,(\varepsilon/\hat{E})^{1/4}$
      \State $q \gets \min\bigl(\max(q, q_{\min}),\, q_{\max}\bigr)$ \Comment{limit how fast $h$ changes}
      \State $h \gets q\,h$ \Comment{new step for retry or next step}
    \EndWhile
  \end{algorithmic}
\end{algorithm}

(b) The local error of the fourth-order formula behaves like $C h^5$, so the per-unit-step estimate behaves like $\hat{E} \approx C h^4$. If we want the next step $qh$ to satisfy $C (qh)^4 \approx \varepsilon$, then $q^4 \hat{E} \approx \varepsilon$ and $q = (\varepsilon/\hat{E})^{1/4}$. The factor $0.84 \approx 2^{-1/4}$ is a safety margin that aims for half the tolerance, which cuts down on rejected steps.

(c) We used $h = 0.25$ initially, $q_{\min} = 0.1$ and $q_{\max} = 4$. Results:

\begin{center}
  \begin{tabular}{@{}ccccc@{}}
    \toprule
    $\varepsilon$ & Accepted & Rejected & $f$-evaluations & $|w(2) - y(2)|$ \\
    \midrule
    $10^{-4}$ & 5  & 0 & 30  & $9.68\times10^{-5}$ \\
    $10^{-6}$ & 14 & 3 & 102 & $1.85\times10^{-6}$ \\
    $10^{-8}$ & 42 & 3 & 270 & $2.00\times10^{-8}$ \\
    \bottomrule
  \end{tabular}
\end{center}

The final errors are of the same order as $\varepsilon$, which is what one hopes for but not what the method guarantees: the controller bounds local, not global, error. At $\varepsilon = 10^{-6}$ the error slightly exceeds the tolerance because the solution grows like $e^t$ and amplifies earlier local errors. Comparing with Table~\ref{tab:errors}, the adaptive code reaches $2\times10^{-8}$ with 270 evaluations, while fixed-step RK4 needs $h = 0.025$ (320 evaluations) to get $2.8\times10^{-8}$. The saving is modest here because the solution is smooth everywhere; on problems with fast transients the adaptive code wins by much more.
\end{solution}

\end{document}
